%! TEX root = part3.tex
% the main TeX file which is intended to compile, :VimtexReload after adjustment
% :h vimtex-tex-root
\documentclass[../gbr_teknik.tex]{subfiles}
% \includeonlyframes{current}
% \setbeameroption{hide notes} % Only slides
% \setbeameroption{show only notes} % Only notes
\setbeameroption{show notes on second screen=right} % Both

\begin{document}

\begin{frame}[mytitle=center,c,mycolor=digiPH_uniblue,light]\frametitle{Tujuan Pembelajaran}

	\framesubtitle{Setelah mempelajari materi ini, mahasiswa mampu:}
	\begin{itemize}
		\item Membaca dan membuat gambar proyeksi sistem Amerika.
		\item Membaca dan membuat gambar proyeksi sistem Eropa.
		\item Membedakan sistem proyeksi Amerika dan Eropa.
	\end{itemize}
\end{frame}

\begin{frame}[mytitle=center,c,mycolor=digiPH_cream,light]Pandangan dalam Gambar Teknik Mesin
	Dalam industri permesinan:

	\begin{itemize}
		\item Gambar harus mudah dibaca dan diinterpretasikan
		\item Harus memiliki pandangan yang cukup
		\item Tidak boleh:
		      \begin{itemize}
			      \item Terlalu sedikit (menyulitkan interpretasi)
			      \item Terlalu banyak (rumit, tumpang tindih / overlap)
		      \end{itemize}
	\end{itemize}


	\tabitem Jumlah pandangan harus seperlunya tetapi lengkap\end{frame}

\begin{frame}[mytitle=standard,c,mycolor=digiPH_leaf,dark]
	\frametitle{Prinsip Penyajian Pandangan}

	\begin{itemize}
		\item Pandangan depan adalah pandangan utama
		\item Pandangan lain berfungsi sebagai:
		      \begin{itemize}
			      \item Pelengkap
			      \item Penjelas
		      \end{itemize}
	\end{itemize}

	Jika satu pandangan sudah cukup menjelaskan bentuk dan ukuran → tidak perlu menambah pandangan lain.\end{frame}



\begin{frame}[mytitle=center,t,mycolor=digiPH_purple,dark]\frametitle{Ketentuan Memilih Pandangan}

	\begin{itemize}
		\item Jangan menggambar lebih dari yang diperlukan
		\item Pilih pandangan yang menunjukkan bentuk paling baik
		\item Utamakan yang memiliki garis tak terlihat paling sedikit
		\item Pandangan kanan lebih utama dari kiri (kecuali lebih informatif)
		\item Pandangan atas lebih utama dari bawah (kecuali lebih informatif)
		\item Pilih pandangan yang mengisi ruang gambar sebaik mungkin
	\end{itemize}
\end{frame}


\begin{frame}[mytitle=standard,t,mycolor=digiPH_gray,light]\frametitle{Proyeksi dalam Gambar Teknik}

	Pandangan divisualisasikan dengan proyeksi lurus

	Terdapat dua sistem:

	\begin{itemize}
		\item Third Angle Projection \textit{(Proyeksi Sistem Amerika / Sudut Ketiga)}
		\item First Angle Projection \textit{(Proyeksi Sistem Eropa / Sudut Pertama)}
	\end{itemize}\end{frame}


\end{document}
